%% Finite syntactic translations; no semantic-completeness inference.

\section{Adjoint presentations of contextual first-order logic}

\begin{annexitem}{annex:p10.adjoint.1}{thm:p10-contextual-predicate-transport}{Discharge, universal closure, and re-instantiation}
\begin{proof}
Write \(\Phi=(\phi_1,\ldots,\phi_m)\) and
\(\Gamma=(x_1:\sigma_1,\ldots,x_n:\sigma_n)\).
Starting from the given finite derivation of
\(\Phi\vdash\psi\,[\Gamma]\), discharge its last hypothesis
by implication introduction, then its next-to-last hypothesis,
and so on. The resulting conclusion is
\[
 \varnothing\vdash D\,[\Gamma],\qquad
 D=\phi_1\Rightarrow
       (\phi_2\Rightarrow\cdots(\phi_m\Rightarrow\psi)\cdots).
\]
When \(m=0\), put \(D=\psi\) and make no implication-introduction
steps. This construction works for lists with repeated formulas:
each application discharges the displayed last occurrence.

Now apply universal introduction first to \(x_n\), then to
\(x_{n-1}\), and continue to \(x_1\).
At every step the removed declaration is last in the current
variable context and no hypothesis remains, so the eigenvariable
condition holds. More explicitly, set \(Q_n=D\) and, for
\(i=n,n-1,\ldots,1\), put
\[
                         Q_{i-1}=\forall_{\sigma_i}((x_i)Q_i).
\]
The successive judgments are
\(\varnothing\vdash Q_i\,[x_1:\sigma_1,\ldots,x_i:\sigma_i]\),
ending in the closed sequent
\(\varnothing\vdash Q_0\,[()]\).
For \(n=0\), this simply means \(Q_0=D\).

Alpha-rename the newly bound variables, and any other bound
variables if necessary, outside the finite set of names in
\(\Delta\) and in all terms \(t_1,\ldots,t_n\).
This changes no formula under the stipulated bound-alpha convention.
Apply context Weaken to obtain
\(\varnothing\vdash Q_0\,[\Delta]\).
Eliminate the outer universal quantifier with
\(t_1:\sigma_1\,[\Delta]\), then the next with
\(t_2:\sigma_2\,[\Delta]\), and continue through \(t_n\).
Each term is well sorted in the retained context \(\Delta\).
Since the bound variables were freshened outside all the terms,
these successive instances give exactly simultaneous substitution,
and the result is
\[
             \varnothing\vdash D[\mathbf t/\mathbf x]\,[\Delta].
\]
This identification is syntactic: in the displayed nested
quantifiers, no substitution for an earlier argument changes a
later substituted term, and no binder captures a free name of it.

For each \(i\), Id gives
\(\phi_i[\mathbf t/\mathbf x]\vdash
  \phi_i[\mathbf t/\mathbf x]\,[\Delta]\).
Apply implication elimination successively to these judgments
and the preceding closed-hypothesis sequent. After \(i\) such
steps the hypothesis list is precisely
\(\phi_1[\mathbf t/\mathbf x],\ldots,
  \phi_i[\mathbf t/\mathbf x]\).
After \(m\) steps the conclusion is the required
\(\Phi[\mathbf t/\mathbf x]\vdash
  \psi[\mathbf t/\mathbf x]\,[\Delta]\).
If \(m=0\), there are no elimination steps.

The original derivation is a subtree of the constructed one.
Every theory-axiom leaf therefore remains in exactly its original
context and form; no renamed or substituted axiom is silently
added to \(\mathcal T\). Only finitely many new rule applications
are made, even if \(\mathcal T\) is infinite.
Choosing the \(t_i\) to be the corresponding renamed variables
proves renaming. Taking \(\Delta\) to be a permutation of
\(\Gamma\) and \(t_i=x_i\) proves context permutation.
When \(\Delta\) is empty, the hypothesis requires
closed terms for every declaration being instantiated.

Only predicate sequents have been transformed. Implication and
universal quantification act on formulas, and the proof contains
no inference from formula equality to an external term equation.
\end{proof}
\end{annexitem}

\begin{annexitem}{annex:p10.adjoint.2}{thm:p10-contextual-adjoint-equivalence}{Every direction of the adjoint and natural-deduction rules}
\begin{proof}
Fix the same arbitrary theory \(\mathcal T\) on both sides.
We give finite derivations of every logical rule from the other
presentation. The term-equation and structural rules are common.

\paragraph{Structural manipulations used below.}
Assumption weakening is derivable: from
\(\Phi\vdash\chi\,[\Gamma]\) and Id
\(\Sigma,\chi\vdash\chi\,[\Gamma]\), Cut gives
\(\Phi,\Sigma\vdash\chi\,[\Gamma]\).
Exchange permits insertion elsewhere in the hypothesis list.
To contract two copies of a finite list, use Exchange to bring
matching occurrences together and Contract to remove the duplicates.
Both operations also apply when a list is empty.

We use Subst only with its literal shared-prefix contexts.
For example, to instantiate a derivation in \(\Gamma,x:\sigma\)
with \(M:\sigma\,[\Gamma]\), its other premise is Eq-Refl
\(M=M:\sigma\,[\Gamma]\); the conclusion is in \(\Gamma\).
From a context \(\Gamma,x:\sigma,y:\tau\), terms \(M,N\)
already well sorted in \(\Gamma\) can be substituted in order:
first \(x:=M\), retaining the suffix \(y:\tau\), then \(y:=N\).
Extending a variable context uses Weaken; quantifier binders are
alpha-renamed fresh before such extension when necessary.

\paragraph{The four propositional bi-rules from natural deduction.}
For \(\&\)-Adj-down, the two premises
\(\Phi\vdash\psi\) and \(\Phi\vdash\psi'\) give
\(\Phi,\Phi\vdash\psi\mathbin{\&}\psi'\) by
\(\&\)-Intr; contract the repeated list.
The two upward directions are exactly \(\&\)-Elim1 and
\(\&\)-Elim2.

For \(\vee\)-Adj-down, use its two premises and Id
\(\psi\vee\psi'\vdash\psi\vee\psi'\) as the three premises
of \(\vee\)-Elim. The conclusion is
\(\Phi,\Phi,\psi\vee\psi'\vdash\theta\).
Contract \(\Phi\) to obtain the required lower sequent.
For the first upward direction, Id and \(\vee\)-Intr1 give
\(\psi\vdash\psi\vee\psi'\). Cut this against
\(\Phi,\psi\vee\psi'\vdash\theta\), then exchange, obtaining
\(\Phi,\psi\vdash\theta\).
For the second direction, Id and \(\vee\)-Intr2 give
\(\psi'\vdash\psi\vee\psi'\); the same Cut gives
\(\Phi,\psi'\vdash\theta\).

The downward implication direction is \(\Rightarrow\)-Intr.
For its upward direction, combine
\(\Phi\vdash\psi\Rightarrow\psi'\) with Id
\(\psi\vdash\psi\) by \(\Rightarrow\)-Elim. This gives
\(\Phi,\psi\vdash\psi'\).
Finally Id gives \(\Phi,\mathsf{false}\vdash\mathsf{false}\);
\(\mathsf{false}\)-Elim gives
\(\Phi,\mathsf{false}\vdash\psi\), the only substantive
direction of \(\mathsf{false}\)-Adj.
All sequents in this paragraph have the same retained context.

\paragraph{The natural-deduction propositional rules from the bi-rules.}
For conjunction introduction with separate antecedents
\(\Phi\vdash\psi\) and \(\Phi'\vdash\psi'\), add hypotheses
and exchange so both antecedents are \(\Phi,\Phi'\).
Apply \(\&\)-Adj-down. Its two upward directions already are
the two conjunction eliminations.

Id \(\psi\vee\psi'\vdash\psi\vee\psi'\) is a lower
instance of \(\vee\)-Adj with empty retained list.
Its upward directions give
\(\psi\vdash\psi\vee\psi'\) and
\(\psi'\vdash\psi\vee\psi'\).
Cut with an arbitrary derivation of either disjunct to obtain
the corresponding natural-deduction introduction.
For elimination, start with its branch premises
\(\Phi,\psi\vdash\theta\) and
\(\Phi',\psi'\vdash\theta\).
Add hypotheses and exchange so that both retained lists are
\(\Sigma=\Phi,\Phi'\). Then \(\vee\)-Adj-down gives
\(\Sigma,\psi\vee\psi'\vdash\theta\).
Cut with the third premise \(\Phi''\vdash\psi\vee\psi'\)
and exchange to give \(\Phi,\Phi',\Phi''\vdash\theta\).

Implication introduction is its downward adjoint direction.
For elimination, the upward direction turns
\(\Phi\vdash\psi\Rightarrow\psi'\) into
\(\Phi,\psi\vdash\psi'\).
Cut with the second premise \(\Phi'\vdash\psi\) and exchange.
For false elimination, use the zero-premise rule to get
\(\mathsf{false}\vdash\psi\), then Cut with
\(\Phi\vdash\mathsf{false}\).
This proves every propositional rule before it is used below.

\paragraph{The quantifier bi-rules from natural deduction.}
The downward universal direction is \(\forall\)-Intr.
For the upward direction, weaken
\(\Phi\vdash\forall_\sigma(\psi)\,[\Gamma]\) to
\(\Gamma,x:\sigma\), freshening the bound argument variable.
Universal elimination with the term \(x:\sigma\,[\Gamma,x:\sigma]\)
gives \(\Phi\vdash\psi(x)\,[\Gamma,x:\sigma]\).

For the downward existential direction, use Id
\(\exists_\sigma(\psi)\vdash\exists_\sigma(\psi)\,[\Gamma]\)
and the supplied upper sequent
\(\Phi,\psi(x)\vdash\theta\,[\Gamma,x:\sigma]\)
as the premises of \(\exists\)-Elim. Exchange its conclusion
into \(\Phi,\exists_\sigma(\psi)\vdash\theta\,[\Gamma]\).
For the upward direction, Id followed by \(\exists\)-Intr
with the well-sorted term \(x\) gives
\[
 \psi(x)\vdash\exists_\sigma(\psi)\,[\Gamma,x:\sigma].
\]
Weaken the supplied lower sequent to this same context and Cut
with the displayed derivation. Exchange yields
\(\Phi,\psi(x)\vdash\theta\,[\Gamma,x:\sigma]\).
The absence of \(x\) from \(\Phi\), and also from \(\theta\)
in the existential bi-rule, is exactly the sheet's side condition.

\paragraph{Natural-deduction quantifier rules from the bi-rules.}
Universal introduction is \(\forall\)-Adj-down.
For elimination, \(\forall\)-Adj-up gives
\(\Phi\vdash\psi(x)\,[\Gamma,x:\sigma]\).
Instantiate \(x\) with the supplied term
\(M:\sigma\,[\Gamma]\), using Eq-Refl and Subst.
The conclusion is \(\Phi\vdash\psi(M)\,[\Gamma]\).

For existential introduction, apply \(\exists\)-Adj-up to Id
\(\exists_\sigma(\psi)\vdash\exists_\sigma(\psi)\,[\Gamma]\),
with empty retained hypothesis list. It yields
\(\psi(x)\vdash\exists_\sigma(\psi)\,[\Gamma,x:\sigma]\).
Instantiate \(x:=M\) with the supplied well-sorted term.
Cut the result against \(\Phi\vdash\psi(M)\,[\Gamma]\).
For existential elimination, apply \(\exists\)-Adj-down to its
second premise \(\Phi',\psi(x)\vdash\theta\,[\Gamma,x:\sigma]\).
The result is \(\Phi',\exists_\sigma(\psi)\vdash\theta\,[\Gamma]\).
Cut with the first premise \(\Phi\vdash\exists_\sigma(\psi)\,[\Gamma]\)
and exchange. No witness is chosen without a well-sorted term,
and no inhabitance assumption has been introduced.

\paragraph{Full equality adjunction from natural deduction.}
First derive a replacement sequent while retaining an arbitrary
parameter context \(K\).
Let \(M,N:\sigma\,[K]\), and let \(\rho\) be a unary formula
abstraction with body in \(K,z:\sigma\), where \(z\) is fresh.
Id followed by implication introduction gives
\[
       \varnothing\vdash\rho(z)\Rightarrow\rho(z)\,[K,z:\sigma].
\]
Apply the registered equality elimination rule using the binary
abstraction \(\chi(u,v)=\rho(u)\Rightarrow\rho(v)\), this
diagonal sequent, and Id
\((M=_\sigma N)\vdash(M=_\sigma N)\,[K]\).
Its conclusion is
\[
          (M=_\sigma N)\vdash\rho(M)\Rightarrow\rho(N)\,[K].
\]
The diagonal premise has no hypotheses, so it satisfies the
actual equality rule's independence condition.

Now suppose the upper sequent of full \(=\)-Adj is given:
\(\Phi\vdash\psi(x,x)\,[\Gamma,x:\sigma]\), where \(\Phi\)
may depend on \(x\). Extend to
\(K=\Gamma,x:\sigma,x':\sigma\).
In the replacement derivation choose \(M=x\), \(N=x'\), and
\(\rho(v)=\psi(x,v)\), with the first \(x\) a retained parameter.
This gives
\[
 (x=_\sigma x')\vdash
            \psi(x,x)\Rightarrow\psi(x,x')\,[K].
\]
Combine it by implication elimination with the weakened given
sequent and exchange the hypothesis lists. The conclusion is
\(\Phi,(x=_\sigma x')\vdash\psi(x,x')\,[K]\), including the
full possible dependence of \(\Phi\) on \(x\).
The fresh variable used in the replacement derivation is appended
after \(K\); no permutation of the original variable context
is needed.

For the upward equality direction, substitute \(x':=x\) in
the given lower sequent, using Eq-Refl in the prefix
\(\Gamma,x:\sigma\). This yields
\(\Phi,(x=_\sigma x)\vdash\psi(x,x)\,[\Gamma,x:\sigma]\).
Equality introduction gives
\(\varnothing\vdash(x=_\sigma x)\,[\Gamma,x:\sigma]\).
Cut removes this equality assumption and leaves the required
upper sequent. The first declaration \(x\) remains present.

\paragraph{Natural-deduction equality from full equality adjunction.}
Id gives
\[
 (x=_\sigma x')\vdash(x=_\sigma x')\,[\Gamma,x:\sigma,x':\sigma].
\]
Use \(=\)-Adj-up with empty retained hypothesis list and
\(\psi(u,v)=(u=_\sigma v)\).
It gives \(\varnothing\vdash(x=_\sigma x)\,[\Gamma,x:\sigma]\).
For any \(M:\sigma\,[\Gamma]\), instantiate \(x:=M\)
by Eq-Refl and Subst, then add the arbitrary hypothesis list.
This is exactly equality introduction.

For equality elimination, its second premise has the form
\(\Phi'\vdash\psi(z,z)\,[\Gamma,z:\sigma]\), with every
member of \(\Phi'\) well formed in \(\Gamma\).
The downward equality direction gives
\[
 \Phi',(z=_\sigma w)\vdash\psi(z,w)\,[\Gamma,z:\sigma,w:\sigma].
\]
Substitute \(z:=M\), retaining \(w\), and then \(w:=N\),
where the supplied terms \(M,N\) are well sorted in \(\Gamma\).
This yields \(\Phi',(M=_\sigma N)\vdash\psi(M,N)\,[\Gamma]\).
Cut with the first premise
\(\Phi\vdash(M=_\sigma N)\,[\Gamma]\), then exchange,
to obtain the registered conclusion.

Every direction has now been supplied by a finite derivation.
Induction on a proof tree replaces each noncommon rule instance
by the corresponding derivation and preserves all theory-axiom
leaves. This proves equivalence for predicate sequents.
The external equational rules and their axiom leaves are identical
in the two presentations, and neither has a rule turning a
predicate sequent into an external equation; hence their external
equation derivability also agrees.
Empty lists, empty ambient contexts, and nullary signature symbols
require no additional rule. The proof uses Cut explicitly and
does not identify the resulting proof trees.
\end{proof}
\end{annexitem}

\begin{annexitem}{annex:p10.adjoint.3}{thm:p10-contextual-equality-reduction}{The parameter-free equality rule and its full parameterization}
\begin{proof}
Retain the common basis and the natural-deduction rules for
\(\Rightarrow\) and \(\forall\).
First, both directions of \(\mathrm{Eq0}\) are literal instances
of full \(=\)-Adj with empty \(\Gamma\) and empty \(\Phi\).
We prove the converse using only the indicated extra logical rules.

\paragraph{Recovering the downward parameterized direction.}
Suppose its upper sequent
\(\Phi\vdash\psi(x,x)\,[\Gamma,x:\sigma]\) is given.
Write \(\Phi=(\phi_1(x),\ldots,\phi_m(x))\), permitting
the formulas to depend on all the variables of \(\Gamma\) as well.
Use implication introduction to discharge these assumptions in
reverse order. With the parameter variables of \(\Gamma\)
temporarily left implicit, put
\[
 D(u,v)=\phi_1(u)\Rightarrow
       \bigl(\phi_2(u)\Rightarrow\cdots
          (\phi_m(u)\Rightarrow\psi(u,v))\cdots\bigr).
\]
When \(m=0\), put \(D(u,v)=\psi(u,v)\).
The result of discharge is
\(\varnothing\vdash D(x,x)\,[\Gamma,x:\sigma]\).
Theorem~\ref{thm:p10-contextual-predicate-transport}, whose proof
uses only the common basis, implication, and universal
quantification, puts this sequent in the order \(x:\sigma,\Gamma\).
This transport remains relative to the same theory axioms.

Write \(\Gamma=(y_1:\tau_1,\ldots,y_k:\tau_k)\).
Universally close these declarations, starting with \(y_k\).
The hypothesis list is empty and each declaration is last at
its step, so universal introduction applies.
Define \(D_k(u,v)=D(u,v)\) and, for
\(j=k,k-1,\ldots,1\), put
\[
             D_{j-1}(u,v)=\forall_{\tau_j}((y_j)D_j(u,v)).
\]
Let \(\chi\) be the binary abstraction with body \(D_0(u,v)\).
For \(k=0\), this means \(\chi(u,v)=D(u,v)\).
All parameters \(y_i\) are now bound, and the derived sequent is
\[
                 \varnothing\vdash\chi(x,x)\,[x:\sigma].
\]
Its formula abstraction has no free parameters besides its two
argument places. Thus Eq0-down applies, giving
\[
            (x=_\sigma x')\vdash\chi(x,x')\,[x:\sigma,x':\sigma].
\]

Freshen the bound \(y_i\) names in this last formula before
appending the context \(\Gamma\) by Weaken.
Then eliminate its \(k\) universal quantifiers using, in order,
the well-sorted terms \(y_i:\tau_i\,[x:\sigma,x':\sigma,\Gamma]\).
The result is
\( (x=_\sigma x')\vdash D(x,x')\,[x:\sigma,x':\sigma,\Gamma]\).
Eliminate the \(m\) implications with Id for each
\(\phi_i(x)\), accumulating the hypothesis list
\((x=_\sigma x'),\Phi\).
Exchange hypotheses and use predicate-context transport to restore
the order \(\Gamma,x:\sigma,x':\sigma\).
The conclusion is exactly
\[
 \Phi,(x=_\sigma x')\vdash
       \psi(x,x')\,[\Gamma,x:\sigma,x':\sigma].
\]
In particular, the entire possible dependence of \(\Phi\) on
\(x\) has been retained inside the discharged formula.

\paragraph{Recovering the upward parameterized direction.}
Id gives
\((x=_\sigma x')\vdash(x=_\sigma x')\,[x:\sigma,x':\sigma]\).
Apply Eq0-up with \(\psi(u,v)=(u=_\sigma v)\).
It gives \(\varnothing\vdash(x=_\sigma x)\,[x:\sigma]\).
Predicate-context transport now gives this reflexivity formula
in \(\Gamma,x:\sigma\): take target context
\(\Gamma,x:\sigma\) and substitute its variable \(x\)
for the single source variable.

Given the lower sequent of full \(=\)-Adj, use the common
Subst rule with Eq-Refl \(x=x:\sigma\,[\Gamma,x:\sigma]\)
to replace the final variable \(x'\) by \(x\).
The result is
\(\Phi,(x=_\sigma x)\vdash\psi(x,x)\,[\Gamma,x:\sigma]\).
Cut with the contextual reflexivity formula just proved.
This recovers the upper sequent with its original retained list.

The argument omits implication steps when \(\Phi\) is empty
and universal steps when \(\Gamma\) is empty. It neither assumes
an inhabitant of any sort nor replaces an axiom by an unproved
renamed instance. Each of the two equality presentations can
therefore replace the other in any finite derivation over the
stated common rules, while all other rules and all axiom leaves
are held fixed. Only implication and universal quantification
were used as extra logical operations; external term equations
entered only through their common reflexivity and substitution
rules, and no equality-reflection rule was used.
\end{proof}
\end{annexitem}
